\chapter{Instrukcja uruchomienia aplikacji REPAIRO}
\label{cha:instrukcja_uruchomienia}

\subsection{Wymagania środowiskowe}

Do poprawnego uruchomienia aplikacji REPAIRO wymagane jest przygotowanie środowiska zgodnie z poniższymi specyfikacjami programowymi:

\begin{itemize}
    \item System operacyjny: \textbf{Windows 10} lub nowszy
    \item Środowisko JDK: \textbf{OpenJDK 23.0.2}
    \item IDE: \textbf{IntelliJ IDEA Community Edition 2024.3.3}
    \item Sterownik JDBC: \textbf{mysql-connector-java-8.0.33} dołączony do repozytorium w katalogu \texttt{mysql-connector-j-9.3.0}.
    \item Baza danych: \textbf{MariaDB} (środowisko \textbf{phpMyAdmin 5.2.1})
    \item Panel sterowania: \textbf{XAMPP Control Panel}
\end{itemize}

\noindent
\textbf{Uwaga:} Domyślnym użytkownikiem serwera bazy danych jest \texttt{root}, a hasło puste (\texttt{""}), zgodnie z domyślną konfiguracją XAMPP.

\subsection{Baza danych}

\begin{itemize}
    \item Do repozytorium dołączony jest plik \texttt{warsztat.sql}, który zawiera strukturę i przykładowe dane bazy danych.
    \item Po uruchomieniu serwera MySQL należy zaimportować bazę danych do środowiska \texttt{phpMyAdmin}.
    \item Domyślny użytkownik: \texttt{root}, Hasło: \texttt{(puste lub zgodne z konfiguracją lokalną)}
    \item Nazwa bazy danych: \texttt{warsztat}
\end{itemize}

\subsection{Konfiguracja połączenia z bazą danych}

W katalogu głównym projektu znajduje się plik:

\begin{center}
\texttt{config-template.properties}
\end{center}

Zawiera on domyślną konfigurację połączenia z bazą danych i \textbf{musi zostać skopiowany oraz dostosowany} przed pierwszym uruchomieniem aplikacji.

\subsubsection*{Kroki:}
\begin{enumerate}
  \item Skopiuj plik:
  \begin{quote}
    \texttt{config-template.properties} \(\rightarrow\) \texttt{config.properties}
  \end{quote}
  \item Otwórz \texttt{config.properties} w edytorze tekstu.
  \newpage\item Uzupełnij dane zgodnie z lokalną konfiguracją użytkownika bazy danych:
\begin{verbatim}
    db.url=jdbc:mysql://localhost:3306/warsztat
    db.user=root
    db.password=
\end{verbatim}
\end{enumerate}

\noindent
\textbf{Uwaga:} Klasa \texttt{DatabaseConnector} pobiera dane z tego pliku w celu nawiązania połączenia z bazą. Dzięki temu nie trzeba edytować kodu źródłowego przy każdej zmianie konfiguracji.

\subsection{Kroki uruchomienia aplikacji REPAIRO}

\begin{enumerate}
  \item \textbf{Klonowanie projektu z GitHub}
  \begin{enumerate}
    \item Przejdź do repozytorium projektu:\\ \url{https://github.com/vLisek/SystemZarzadzaniaSerwisemSamochodowym.git}
    \item Pobierz repozytorium lokalnie:
    \item {\small\texttt{git clone https://github.com/vLisek/SystemZarzadzaniaSerwisemSamochodowym.git}}
    \item Otwórz folder projektu w IntelliJ IDEA Community Edition.
  \end{enumerate}

  \item \textbf{Dodanie sterownika JDBC}
  \begin{enumerate}
    \item Umieść folder \texttt{mysql-connector-j-9.3.0} w katalogu projektu
    \item Dodaj go jako bibliotekę w: \texttt{File > Project Structure > Modules > Dependencies}
  \end{enumerate}

  \item \textbf{Import bazy danych}
  \begin{enumerate}
    \item Uruchom \texttt{XAMPP Control Panel} i włącz moduł MySQL
    \item Otwórz \texttt{phpMyAdmin} w przeglądarce: \url{http://localhost/phpmyadmin}
    \item Utwórz nową bazę danych: \texttt{warsztat}
    \item Zaimportuj plik \texttt{warsztat.sql}
  \end{enumerate}

  \item \textbf{Uruchomienie aplikacji}
  \begin{enumerate}
    \item Kliknij przycisk \texttt{Run} w klasie \texttt{Main.java}
    \item Aplikacja otworzy ekran logowania
  \end{enumerate}
\end{enumerate}

\newpage
\subsection{Dane testowe (logowanie)}

\begin{itemize}
  \item \textbf{Administator}
  \begin{itemize}
    \item Nazwa użytkownika: \texttt{admin}
    \item Hasło: \texttt{admin}
  \end{itemize}

  \item \textbf{Mechanik}
  \begin{itemize}
    \item Nazwa użytkownika: \texttt{kszymanski}
    \item Hasło: \texttt{1234}
  \end{itemize}
\end{itemize}

\noindent
Po wykonaniu powyższych kroków środowisko jest gotowe do pracy z aplikacją REPAIRO.
